% !TeX program = pdflatex
% conecthus/pipeline.tex — especificação da IR Corpo/Conecthus (v0.2).
% Promovida a eixo da arquitetura: corpus/docs/arquitetura.tex §sec:pipeline.
% Auto-contido: pdflatex pipeline.tex
% Fonte: decisão do coordenador + fundamento.tex + papers/corpo\_topologico.tex + arquitetura.tex.
% Isto NÃO é ainda a linguagem executável --- é a especificação dela.
\documentclass[11pt,a4paper]{article}
\usepackage[utf8]{inputenc}\usepackage[T1]{fontenc}\usepackage[portuguese]{babel}
\usepackage{amsmath,amssymb,amsthm}
\usepackage[margin=2.3cm]{geometry}
\usepackage{booktabs}\usepackage{xcolor}
\usepackage[hidelinks]{hyperref}

\definecolor{cinza}{gray}{0.35}
\newcommand{\code}[1]{\texttt{\small #1}}
\newcommand{\medido}[1]{\par\medskip\noindent{\small\color{cinza}\textbf{Medido.} #1}\par\medskip}
\newtheorem{teorema}{Teorema}
\newtheorem{definicao}[teorema]{Definição}
\newtheorem{proposicao}[teorema]{Proposição}
\newtheorem{obrigacao}[teorema]{Obrigação}
\newtheorem{observacao}[teorema]{Observação}
\newcommand{\dg}{^{\dagger}}

\title{\textbf{O pipeline}\\[3pt]
  {\large especificação v1.6 --- portão da Pátria (local contra o ar)}}
\author{Aarão Melo Lopes}
\date{13 de agosto de 2026}

\begin{document}
\maketitle

\begin{abstract}
\noindent O pipeline chega ao ar pela mesma régua: local (corpo tem
\code{pipeline.tex}) contra live (HTTP 200). $0\wedge0$ não fecha;
afirmar no ar sem fetch (réu) não fecha.
\code{deploy.claim}; chip \code{D}; \code{tools/patria.sh}.
Claim\,$\neq$\,Result.
\end{abstract}

\tableofcontents

%----------------------------------------------------------------------
\section{O que a v0.1 acertou, e o que o gerente corrige}\label{sec:v01-v02}

A v0.1 formalizou \textsf{Claim}, \textsf{Step}/\textsf{Back}, $R$, a gramática
do fundamento e o contrato «nenhum PASS nu». Mantém-se. Corrige-se:

\begin{center}\small
\begin{tabular}{@{}p{5.2cm}p{7.0cm}@{}}
\toprule
v0.1 & v0.2 \\
\midrule
\code{residual} dentro do Claim como expressão & Claim = especificação; Result = observação \\
Controlo na sintaxe v0.1 & Controlo = runtime \emph{acima} da IR \\
árvore larga (gut/pdca/ishikawa/\ldots) & três claims só; o resto são instâncias \\
paper como par da IR & \code{*.claim} é fonte; \code{.tex} é backend \\
pipeline = linguagem & pipeline = especificação; IR executável vem a seguir \\
\bottomrule
\end{tabular}
\end{center}

\begin{obrigacao}[Claim não escreve o seu próprio PASS]\label{ob:claim-result}
Nenhuma projeção --- nem o Claim --- declara o resultado. Resultado só existe
depois de $\mathrm{STEP}\to\mathrm{BACK}\to\mathrm{MEASURE}\to R$. Escrever
\code{residual = 0} no Claim e acreditar nele é exactamente o problema que o
fundamento elimina~\cite{fundamento}.
\end{obrigacao}

%----------------------------------------------------------------------
\section{Enquadramento: as quatro camadas do Corpo de Peano}\label{sec:peano}

A Prop.~\emph{central-operacional}~\cite{peano} dualiza o Teorema Central em
quatro camadas sem as colar. A IR mapeia nelas:

\begin{center}\small
\begin{tabular}{@{}lll@{}}
\toprule
camada Peano & objecto na torre & na IR \\
\midrule
estrutura & retração $\pi_k$ & tipagem / domínio do Claim \\
dinâmica & Maestro (projecta) & \textsf{STEP} \\
métrica & Metrónomo ($\lambda^{+}+\lambda^{-}=0$) & \textsf{BACK}+\textsf{MEASURE}$\to R$ \\
controlo & histerese / rede dual & runtime \emph{acima} (RETAIN/MOVE) \\
música & Conservatório$\to$Pera$\to$Batuta & backends (LaTeX, C, teste, relatório) \\
\bottomrule
\end{tabular}
\end{center}

\[
\boxed{
\begin{array}{c}
\text{TEXTO / IDEIA}\\
\downarrow\\
\mathrm{EXTRACT}\to\mathrm{IR}\\
\downarrow\\
\underbrace{\mathsf{STEP}}_{\text{Maestro}}
\;\to\;
\underbrace{\mathsf{BACK}+\mathsf{MEASURE}}_{\text{Metrónomo}}
\;\to\; R\\
\downarrow\\
R=0\Rightarrow\mathrm{CLOSE}\qquad
R\neq0\Rightarrow\mathrm{REOPEN}\\[4pt]
\text{MUTATE fora do caminho normal}\\[2pt]
\underbrace{\mathrm{CONTROL}}_{\text{histerese}}
\text{ só depois da IR fechar}
\end{array}
}
\]

\begin{observacao}[LLM $\neq$ Maestro]
O modelo de linguagem pode \emph{extrair} Claims (estágio INGEST/EXTRACT). Não
é o Maestro: não projecta $\pi_k$; não atesta $\lambda^{+}+\lambda^{-}=0$. A
atestação é o medidor --- Metrónomo --- fora do modelo~\cite{peano,fundamento}.
\end{observacao}

\begin{observacao}[partitura $\Pi$ e \code{*.claim}]
Na torre, $\Pi$ é o \emph{quê}: transporte, não fonte das leis~\cite{peano}.
Na IR, o ficheiro \code{.claim} desempenha o mesmo papel operacional: especifica
o quê; o executor e o medidor realizam a dinâmica e a métrica. O paper LaTeX é
Conservatório/folha --- realização, não fonte.
\end{observacao}

%----------------------------------------------------------------------
\section{Arquitectura do compilador}\label{sec:arch}

\begin{verbatim}
                 TEXTO / IDEIA
                      |
                      v
                 [ EXTRACT ]
                      |
                      v
              +---------------+
              |   IR (fonte)  |
              | Claim Step    |
              | Back Measure  |
              | Residual*     |
              | Mutation      |
              +-------+-------+
                      |
           +----------+----------+
           v          v          v
         PAPER      TEST       CODE
           |          |          |
           +----------+----------+
                      v
                   EXECUTE          <-- Maestro
                      |
                      v
                     BACK
                      |
                      v
                   MEASURE          <-- Metrónomo
                      |
                      v
                  RESIDUAL R
                   /       \
                R = 0     R != 0
                  |          |
               CLOSE       REOPEN
                  |          |
               CRISTAL   estagio responsavel
\end{verbatim}
\noindent (*) Residual no diagrama é o \emph{cálculo} do runtime, não um campo
que o Claim preenche com o valor desejado.

%----------------------------------------------------------------------
\section{Núcleo mínimo: sete primitivas + duas decisões}\label{sec:nucleo}

\begin{definicao}[primitivas]\label{def:prim}
\[
\{\mathsf{CLAIM},\,\mathsf{STEP},\,\mathsf{BACK},\,\mathsf{MEASURE},\,
\mathsf{RESIDUAL},\,\mathsf{MUTATE},\,\mathsf{CLASSIFY}\}
\]
e as decisões $\{\mathsf{CLOSE},\,\mathsf{REOPEN}\}$.
\end{definicao}

Formalmente o caminho normal:
\[
x
\xrightarrow{\mathrm{STEP}}
x'
\xrightarrow{\mathrm{BACK}}
\hat{x}
\xrightarrow{\mathrm{MEASURE}}
R(x,\hat{x}).
\]
\[
R=0\Rightarrow\mathrm{CLOSE},
\qquad
R\neq0\Rightarrow\mathrm{REOPEN}.
\]

A mutação fica \textbf{fora} do caminho normal:
\[
\mathrm{CLAIM}\to\mathrm{EXECUTE}\to R,
\qquad\text{depois}\qquad
\mathrm{MUTATE}(\mathrm{CLAIM})\to\mathrm{EXECUTE}\to R'.
\]
Isto distingue \emph{fechou} de \emph{sobreviveu à mutação}.

\begin{definicao}[\textsf{Claim} vs.\ \textsf{Result}]\label{def:claim-result}
O Claim declara lei, objecto, operação, volta \emph{como procedimento},
invariante \emph{como predicado}, e mutação. Não declara o valor de $R$.
O Result é produzido pelo runtime:
\begin{verbatim}
result {
    forward  = ...
    reverse  = ...
    residual = ...
    mutation = SURVIVED | BROKEN | UNTESTED
    class    = STRUCTURAL | CONVENTIONAL | REALIZATION
               | HYPOTHESIS | FALSE | INCOMPLETE
}
\end{verbatim}
\end{definicao}

%----------------------------------------------------------------------
\section{Pipeline operacional (estágios 0--12)}\label{sec:estagios}

\begin{center}\small
\begin{tabular}{@{}clp{7.4cm}@{}}
\toprule
 & estágio & acção \\
\midrule
0 & INGEST & texto / paper / conversa / código existente \\
1 & EXTRACT & extrair claims \\
2 & NORMALIZE & prosa $\to$ primitivas IR \\
3 & TYPE & verificar tipos e domínios ($\sim\pi_k$ tipada) \\
4 & GENERATE & gerar backends (paper/teste/código) \\
5 & EXECUTE & executar Step (Maestro) \\
6 & BACK & executar a volta \\
7 & MEASURE & medir (Metrónomo) \\
8 & RESIDUAL & calcular $R$ \\
9 & MUTATE & perturbar Claim \\
10 & CLASSIFY & estrutural / convencional / \ldots \\
11 & CLOSE & cristalizar \\
12 & REOPEN & devolver ao estágio que falhou \\
\bottomrule
\end{tabular}
\end{center}

\begin{proposicao}[REOPEN é depuração epistemológica]\label{prop:reopen}
O estágio~12 não precisa de «saber programar». Encaminha:
\begin{itemize}\itemsep2pt
\item TYPE falhou $\to$ NORMALIZE/TYPE;
\item BACK falhou $\to$ definição do operador;
\item MEASURE falhou $\to$ medidor;
\item MUTATION falhou / não correu $\to$ class = HYPOTHESIS ou UNTESTED.
\end{itemize}
O pipeline deixa de ser só CI: torna-se sistema de depuração da afirmação.
\end{proposicao}

%----------------------------------------------------------------------
\section{Formato textual mínimo (primeiro compilador)}\label{sec:formato}

Não se começa por parser sofisticado. Formato pequeno:

\begin{verbatim}
claim ParetoClosure
law 6
object Pareto
step sum(percentages)
back recompute(percentages)
measure sum(percentages)
invariant sum == 100
mutate numerator += 1
end
\end{verbatim}

\noindent O compilador produz AST canónica \code{Claim(...)} e projecta:
\begin{verbatim}
claim -> JSON | C | Python | LaTeX | test | relatorio
\end{verbatim}

\begin{obrigacao}[fonte de verdade]\label{ob:fonte}
A fonte passa a ser \code{claims/*.claim} (ou \code{corpo/*.ir}). O
\code{.tex} deixa de ser fonte de verdade: é backend gerado ou anotado a
partir da IR. Nenhuma projeção declara o seu próprio resultado.
\end{obrigacao}

\subsection*{Regra de ouro (contrato do projecto)}
\begin{verbatim}
A IR e a fonte de verdade.
Paper, teste, codigo e visualizacao sao projecoes.
Nenhuma projecao pode declarar o seu proprio resultado.
Resultado so existe depois de: STEP -> BACK -> MEASURE -> RESIDUAL
R = 0  -> CLOSE
R != 0 -> REOPEN
MUTATION determina CLASSIFICATION.

PASS sem residual e invalido.
PASS sem back e invalido.
STRUCTURAL sem mutation e hipotese.
MEASURE sem reverse e medida unilateral.
\end{verbatim}

%----------------------------------------------------------------------
\section{Três claims de arranque --- e só três}\label{sec:tres}

Não se começa pelas nove ferramentas Lean. Três comportamentos distintos;
se a IR os expressa sem excepções especiais, começa a ser linguagem.

\subsection*{1. Pareto --- invariância numérica}
\[
R=\Sigma-100\%.
\]
Step, Back, medida e mutação já medidos no fundamento
(\code{tests/lean\_cruz.c})~\cite{fundamento}.

\begin{verbatim}
claim ParetoClosure
law 6
object ParetoTable
step sum(percentages)
back recompute(percentages)
measure sum(percentages)
invariant sum == 100
mutate numerator += 1
classify structural
end
\end{verbatim}
O runtime (não o ficheiro) emite \code{result\{ residual=0, mutation=BROKEN, \ldots\}}.

\subsection*{2. Why --- ponto fixo}
\[
A_{n+1}=\varepsilon(A_n),\qquad R=A_{n+1}-A_n.
\]
O $5$ não é a lei: é \code{default\_depth} (convenção). A estrutura é o
ponto fixo~\cite{fundamento}.

\begin{verbatim}
claim WhyFixedPoint
law 5
object WhyChain
step erode
back next
measure root
invariant next(root) == root
mutate depth = 1..8
classify depth = convention
end
\end{verbatim}

\subsection*{3. Kanban --- mutação estrutural}
\begin{verbatim}
claim KanbanThree
law 3
object Kanban
step transit(states={TODO,PROGRESS,DONE})
back invert_flow
measure collisions
invariant collisions == 0
mutate merge(PROGRESS, DONE)
classify structural
end
\end{verbatim}

\noindent Triplo de exercícios:
\begin{verbatim}
Pareto  -> invariancia numerica
Why     -> ponto fixo
Kanban  -> mutacao estrutural
\end{verbatim}

\section{Instâncias Lean --- GUT e PDCA (v0.8)}\label{sec:instancias}

O gerente~\cite{eval} proíbe árvores \code{gut/}/\code{pdca/} independentes:
primeiro demonstram que são \textbf{só} instâncias da IR. Feito:

\subsection*{GUT --- produto que ordena}
Par do Pareto (soma mede / produto ordena)~\cite{fundamento}.
Escala $G,U,T\in\{1,\ldots,5\}$; mutação \code{allow\_zero} quebra monotonia
(\code{tests/lean\_cruz.c} \S L2).

\begin{verbatim}
claim GUTProduct
law 6
object GUTMatrix
step product(G,U,T)
back recompute(score)
measure score
invariant score == G*U*T && G,U,T in 1..5
mutate allow_zero
classify structural
end
\end{verbatim}

\subsection*{PDCA --- ciclo fechado pela leitura}
O Check precisa de alvo numérico no Plan; sem número o residual é
ilegível~\cite{fundamento}. Mutação \code{drop\_numeric\_target}.

\begin{verbatim}
claim PDCACycle
law 6
object PDCA
step plan_do_check_act
back return_to_plan
measure check_residual
invariant period == 4 && plan.has_numeric_target
mutate drop_numeric_target
classify structural
end
\end{verbatim}

\medido{\code{tests/claim\_ir.c} \S C4: GUT fecha em $(3,4,5)$ com score $60$;
$0$ na escala REOPEN; PDCA fecha com alvo; sem número não fecha; Check lê
$R=\mathrm{plan}-\mathrm{done}$. Wasm ids $3$/$4$ em \code{claim.wasm}.}

\subsection*{5W2H --- dilatação (par do Why)}
Erosão (Why) escolhe a raiz; dilatação (5W2H) escreve acções. Invariante:
raiz $\in$ plano. Mutação \code{drop\_root} quebra.

\begin{verbatim}
claim FiveW2H
law 6
object ActionPlan
step dilate(root)
back erode(actions)
measure containment
invariant root in actions
mutate drop_root
classify structural
end
\end{verbatim}

\subsection*{Ishikawa --- $k$ livre}
Re-partir causas em $k\in\{2,\ldots,12\}$: escolha \emph{por causa} não
move a raiz (resíduo $0$); o $6$ é convenção~\cite{fundamento}.

\begin{verbatim}
claim IshikawaFree
law 6
object CauseMap
step group(causes, k)
back choose_by_cause
measure root_stable
invariant root stable under repartition
mutate k = 2..12
classify k = convention
end
\end{verbatim}

\medido{\code{tests/claim\_ir.c} \S C5: FiveW2H fecha com raiz; Ishikawa
fecha em $k=6$ e $k=3$ com a mesma raiz (mutação SURVIVED); raiz movida
REOPEN. Wasm ids $5$/$6$.}

\subsection*{VSM e Fluxograma --- o dual nomeado}
Os dois membros (actual / futuro) têm de existir; tirar o futuro quebra
(involução pela metade)~\cite{fundamento}.

\begin{verbatim}
claim VSMDual
law 1
object ValueStream
step name_atual_futuro
back involution_swap
measure both_named
invariant atual and futuro present
mutate drop_futuro
classify structural
end

claim FluxogramaDual
law 1
object ProcessMap
step name_nodes_atual_futuro
back involution_swap
measure nodes_paired
invariant every node has atual and futuro
mutate drop_futuro_node
classify structural
end
\end{verbatim}

\medido{\code{tests/claim\_ir.c} \S C6: VSM fecha com dual; sem futuro REOPEN;
Fluxograma fecha com $n$ nós pareados; parcial $R=n-\mathrm{paired}$.
Wasm ids $7$/$8$. Inventário LMS completo na IR.}

\section{Fronteira entre Claims --- Lei~7 (v1.1)}\label{sec:fronteira}

Fundir artefactos num só contexto torna a inconsistência ilegível. A forma
obrigada~\cite{fundamento}: \textbf{dois corpos e a interface}, resíduo
contável --- sem palavra-chave nova na gramática além de mais
\code{.claim}.

\begin{verbatim}
claim FrontierGUT5W2H     # inversões ranking ↔ ordem de ações
claim FrontierWhyIshikawa # raiz Why ∈ mapa Ishikawa
claim FrontierPDCAVSM     # plan.target == vsm.futuro
\end{verbatim}

\medido{\code{tests/claim\_fronteira.c}: par consistente $R=0$; swap
adjacente apanhado na fronteira (interna continua verde); raiz fora /
alvo$\neq$futuro REOPEN. Wasm ids $9$--$11$. Espelha
\code{tests/lean\_controlo.c} \S N2.}

\section{Volta na estação --- projectado contra medido (v1.4)}\label{sec:estacao}

A validação interna lê o artefacto contra si próprio --- metade do par.
A outra metade é o mundo~\cite{fundamento}: lead time projectado contra o
medido na estação. Sem essa coluna, $65\%$ e $200\%$ (estação~04) ficam
projeções. A referência escrita pelo réu (números inferidos) valida-se
sempre --- mutação \code{copy\_reference}.

\begin{verbatim}
claim StationResidual
law 0
object Station
step implant
back measure_station
measure abs(projected - measured)
invariant measured from station not from model
mutate copy_reference
classify structural
end
\end{verbatim}

$R=|P-M|$. Fecha só se a leitura é externa (MES) e $R=0$. Chip \code{E}:
\code{lead:65 estacao:65}.

\medido{\code{tests/claim\_ir.c} \S C7; \code{claim.wasm} id $12$;
\code{tests/estacao\_front.js}.}

\section{O $R$ no banco --- documento, não pesos (v1.5)}\label{sec:banco}

A Obrigação~\cite{fundamento} \emph{ob:banco} proíbe empurrar o fecho do
ciclo para dentro da interface (fine-tune / RLHF sobre a correcção).
O que a lei permite já está no sistema: o corpus recuperável. O desfecho
da estação vira par (\code{fala}, \code{resposta}) ---
\code{emit}$\circ$\code{parse}=\code{id} --- e a recuperação trá-lo ao
próximo prompt. O modelo continua sem estado; o banco lembra.

\begin{verbatim}
claim BancoVolta
law 0
object Bank
step write_document
back retrieve
measure compare
invariant source == retrieved
mutate alter_doc
classify structural
end
\end{verbatim}

Canónico: \code{implante P=p M=m R=|P-M| fonte=MES|reu}.
A estação pode ficar aberta ($R\neq0$) e o banco \emph{fechar} --- o
documento está intacto. Mutação \code{alter\_doc} denuncia.

Chip \code{B}: após a leitura na estação, \code{aprende}; «mostra o banco»
recupera. Sem daemon, o documento ainda é o texto canónico (não um peso).

\medido{\code{tests/claim\_ir.c} \S C8; \code{claim.wasm} id $13$;
\code{tests/estacao\_front.js} \S E4--E5; \code{tools/pipeline.sh}.}

\section{Portão da Pátria --- local contra o ar (v1.6)}\label{sec:patria}

Publicar não é um passo à parte da IR: é a volta contra o sítio.
O projectado é o corpo medido (\code{tools/corpo.sh}:
\code{conecthus/pipeline.tex} na lista). O medido é o GET
\texttt{/corpo/conecthus/pipeline.tex} em
\href{https://goldenkingdom.patriatechnology.com}{goldenkingdom.patriatechnology.com}.
Os dois a zero não fecham --- não há artefacto. A afirmação sem fetch
é a referência do réu.

\begin{verbatim}
claim DeployPatria
law 0
object Publication
step build
back fetch_live
measure compare
invariant local and live
mutate drop_pipeline
classify structural
end
\end{verbatim}

Chip \code{D}: «mostra o deploy». O \code{publica.yml} recusa o
artefacto se o health check não vir o pipeline e o \code{claim.wasm}.
O push em si fica ao dono --- o Claim observa, não dispara o git.

\medido{\code{tests/claim\_ir.c} \S C9; \code{claim.wasm} id $14$;
\code{tests/patria\_front.js}; \code{tools/patria.sh}.}

%----------------------------------------------------------------------
\section{Controlo --- runtime acima da IR (v0.4)}\label{sec:controlo}

O Controlo de Histerese~\cite{peano} (\code{thm:controle-histerese})
\textbf{não} entra no parser. É camada acima, já realizada:

\begin{verbatim}
IR  (Claim, Step, Back, Measure -> Result)
 |
 v
CONTROL   conecthus/core/control.c
 |        \
RETAIN   MOVE / RETRACT
\end{verbatim}

\[
\boxed{
\operatorname{Gerar}
\to
\operatorname{Medir}(D)
\to
\operatorname{Seleccionar}
\to
\{\mathrm{RETAIN},\,\mathrm{MOVE},\,\mathrm{RETRACT}\}
}
\]

\begin{center}\small
\begin{tabular}{@{}ll@{}}
\toprule
acção & quando \\
\midrule
\code{RETAIN} & $R=0$ e $D\preceq\Theta$ --- estado preservado \\
\code{RETRACT} & $R\neq0$ --- REOPEN do pipeline \\
\code{MOVE} & $R=0$ mas $D\npreceq\Theta$ --- fora da borda / muda andar \\
\bottomrule
\end{tabular}
\end{center}

\noindent Dualidade: histerese $=$ retenção; controlo $=$ selecção. O Controlo
já corre acima da IR (\code{conecthus/core/control.c} e
\code{backends/control/runtime.c}$\to$\code{control.wasm}):
$R=0\wedge D\preceq\Theta\Rightarrow\mathrm{RETAIN}$;
$R\neq0\Rightarrow\mathrm{RETRACT}$; $R=0\wedge D\npreceq\Theta\Rightarrow\mathrm{MOVE}$.
Oito eixos em $\mathcal{V}$: L1, $|R|$, $|n|$, mut, \emph{caixa},
\emph{teclado} ($d_K$), \emph{fonético} ($d_{\phi}$), \emph{forma}
(\code{conecthus/core/eixos\_texto.c}; pares reais \texttt{cafe}/\texttt{café}).
Sem palavra-chave \code{control} no \code{.claim}.

\begin{observacao}
Na torre Peano o controlo realiza a atestação na interface --- não substitui
o Teorema Central~\cite{peano}. Misturá-lo na sintaxe do Claim seria colar
camadas.
\end{observacao}

%----------------------------------------------------------------------
\section{Árvore mínima e ordem de construção}\label{sec:arvore}

\begin{verbatim}
conecthus/
  backends/   c ptx wasm llvm glsl haskell dafny isabelle latex claim isa
              manifesto.json   <-- nenhuma privilegiada; C subset -> traduz -> wasm
              latex/claim_to_tex.c  <-- projecao LaTeX (sem Result)
  claims/     9 LMS + pipeline + 3 fronteiras + estacao + banco + deploy
  core/       execute control eixos_texto fronteira
tests/claim_fronteira.c       Lei 7 — resíduo entre Claims
app/src/fronteira.js          chip F na assistente
app/src/estacao.js            chip E — P contra M; documento B
app/src/patria.js             chip D — local contra o ar
tests/estacao_front.js        medidor UI estação / banco
tests/patria_front.js         medidor UI Pátria
tools/pipeline.sh             corpus + fronteira + Controlo + banco + Pátria
tools/patria.sh               portão local (--live = GET)
app/src/fronteira.js          chip F na assistente
tests/fronteira_front.js      medidor UI fronteira
  lang/       claim.h  parse_claim.c   <-- parser (v0.3)
  core/       execute.h/.c   control.h/.c   <-- IR + histerese acima
  fundamento.tex
  pipeline.tex
corpus/docs/arquitetura.tex  §pipeline = terceiro eixo
tools/sobe_backends_wasm.sh
tests/backends_wasm.js        #TOTAL 26 0
tests/claim_ir.c              parser+executor+PipelineClosure
tests/claim_control.c         RETAIN/MOVE/RETRACT + 8 eixos V
assets/.../control.wasm       Controlo na ponte
app/src/controlo.js           chip C + eixos texto
conecthus/core/eixos_texto.c  caixa/teclado/φ/forma (dados reais)
tests/controlo_front.js       medidor UI
corpus/fala/conversa_pipeline.tex  corpus
\end{verbatim}

\noindent A ordem do gerente \textbf{fecha}:
\[
\boxed{
\text{IR}
\to\text{parser}
\to\text{executor}
\to\text{medidor}
\to\text{mutação}
\to\text{10 claims}
\to\text{wasm/LaTeX}
\to\text{Controlo}
}
\]

\medido{\code{tests/claim\_ir.c}: dez \code{.claim} (9 LMS + PipelineClosure).
Controlo: \code{tests/claim\_control.c}.}

\subsection*{O pipeline mede-se a si mesmo}
O primeiro Claim meta pode ser a própria régua:

\begin{verbatim}
claim PipelineClosure
law 0
object Pipeline
step compile
back reconstruct
measure compare
invariant source == reconstructed
mutate alter_ast
end
\end{verbatim}

\noindent Lei~0 do fundamento: todo passo reverte, ou paga --- aplicada ao
compilador que a implementa. No eixo Peano: Maestro compila (projecta);
Metrónomo compara fonte e reconstrução ($R=0$ na volta).

%----------------------------------------------------------------------
\section{O que isto obriga ao LMS e à torre}\label{sec:obriga}

\begin{itemize}\itemsep3pt
\item onde o LMS já reverte (re-deriva GUT/Pareto/PDCA), já é
  EXECUTE+BACK+MEASURE~\cite{fundamento};
\item PASS sem back / sem residual é inválido --- limitações 4.5.x onde a
  volta falta;
\item STRUCTURAL sem mutation é hipótese --- a fronteira
  Kanban/Ishikawa/Porquês já medida por mutação;
\item MEASURE sem reverse é medida unilateral --- Observação da validação
  parcial (artefacto contra si vs.\ contra o mundo);
\item Controlo (acaso ao lado do ganho) entra \emph{depois} dos três claims
  fecharem;
\item na torre: não identificar rede dual com Gentil--Hurwitz; não
  identificar backend musical com a fonte das leis~\cite{peano}.
\end{itemize}

\[
\boxed{
\mathrm{Law}\to\mathrm{Claim}\to\mathrm{Step}\to\mathrm{Back}
\to\mathrm{Measure}\to R\to\mathrm{Mutate}\to\mathrm{Classify}
\to\{\mathrm{Close},\mathrm{Reopen}\}
}
\]
\noindent A especificação v1.6 fecha o portão da Pátria (corpo contra
HTTP; $0\wedge0$ recusado). A entrega seguinte é o push --- o dono
dispara; o Claim só observa.

\begin{thebibliography}{9}
\bibitem{arquitetura} A.~M.~Lopes. \emph{A arquitetura}
  (\code{corpus/docs/arquitetura.tex}) --- três eixos: porta, torre, pipeline
  (\S\,pipeline); «todo passo reverte, ou paga».
\bibitem{eval} Coordenador. Direcionamento arquitectónico (08/2026).
  (13~ago~2026) --- Claim\,$\neq$\,Result; Controlo acima; três claims;
  ordem IR$\to$parser$\to$executor.
\bibitem{fundamento} A.~M.~Lopes. \emph{O Lean reverte}
  (\code{conecthus/fundamento.tex}).
\bibitem{peano} A.~M.~Lopes. \emph{Corpo de Peano}
  (\code{papers/corpo\_topologico.tex}) --- $\pi_k$, Maestro/Metrónomo,
  Prop.~central-operacional, histerese, rede dual.
\bibitem{lms} D.~N.~Lopes et al. \emph{LLM-Driven Lean Manufacturing System}.
  JCSTS 8(8): 91--102, 2026.
\bibitem{teoria} A.~M.~Lopes. \emph{A teoria} (\code{teoria.tex}).
\end{thebibliography}

\vfill
{\footnotesize\color{cinza}\noindent
Especificação v0.2 --- decisão do coordenador enquadrada no Corpo de Peano.\\
\href{https://goldenkingdom.patriatechnology.com}{goldenkingdom.patriatechnology.com}.
CC BY 4.0 (textos) $\cdot$ MIT (código).\par}

\end{document}
